% ----------------------------------------------------------
% Capítulo 1 — Introdução
% ----------------------------------------------------------
\chapter{Introdução}
\label{cap:introducao}

% ---
\section{Contextualização}
\label{sec:contextualizacao}
% ---

A informação geoespacial tem papel estratégico na defesa, no planejamento territorial, na gestão ambiental e na tomada de decisão do Estado brasileiro. O Decreto n.º 6.666/2008 instituiu a Infraestrutura Nacional de Dados Espaciais (INDE), que organiza o acesso integrado aos dados geoespaciais produzidos pelos órgãos federais e inclui a disseminação e o uso desses dados entre os seus objetivos \cite{inde2010decreto}. No Exército Brasileiro, a Diretoria de Serviço Geográfico (DSG) orienta as atividades referentes a imagens e informações geográficas, com vistas à produção cartográfica, ao estabelecimento de normas e ao suprimento do material a seu cargo \cite{dsg2024missao}. A DSG mantém o Banco de Dados Geográficos do Exército (BDGEx), criado para disponibilizar na Internet os dados geoespaciais da instituição, conforme o decreto da INDE \cite{xavier2014bdgex}, e por meio dele o acervo cartográfico oficial fica acessível ao público militar e à sociedade civil \cite{bdgex2024}. Cinco Centros de Geoinformação (CGEO), sediados em Porto Alegre, Brasília, Olinda, Manaus e Rio de Janeiro, produzem esse acervo segundo a Especificação Técnica para Produtos de Conjuntos de Dados Geoespaciais (ET-PCDG) \cite{etpcdg2016} e o Perfil de Metadados Geoespaciais do Brasil (Perfil MGB) \cite{perfilmgb2009}, e cada CGEO atende a uma região, a sua área de suprimento. Os produtos são catalogados em serviços compatíveis com os padrões do \textit{Open Geospatial Consortium} (OGC), em particular o \textit{Catalogue Service for the Web} (CSW) \cite{ogccsw2007}, e o catálogo organiza os metadados num modelo de dados rígido.

A busca nesse catálogo exige que o usuário conheça as convenções técnicas que organizam o acervo. Ele precisa identificar a folha, unidade da articulação sistemática das cartas, pelo Índice de Nomenclatura (INOM), código que subdivide a folha da Carta Internacional do Mundo ao Milionésimo até a escala pedida, ou pelo Mapa Índice (MI), número da folha 1:100.000 que as folhas 1:50.000 e 1:25.000 herdam com sufixos \cite{ibge1999nocoes}. Precisa também saber quais escalas o catálogo oferece, conhecer os tipos de produto do Sistema Cartográfico Nacional (SCN) \cite{scn2016}, como ``SCN Carta Topográfica Matricial'', e reconhecer o CGEO responsável por cada folha. Esses conceitos são detalhados na Seção~\ref{sec:res-modelagem}, junto das ferramentas que os codificam. A interface de busca do catálogo compara o texto digitado com os campos do banco, sem interpretar a linguagem natural, de modo que só encontra o produto quem escreve os termos técnicos do acervo. A Tabela~\ref{tab:consulta-tecnica} compara esse modo de consulta com o modo desejado.

\begin{fixo}{table}
\centering
\caption{Consulta técnica e consulta em linguagem natural}
\label{tab:consulta-tecnica}
\begin{tabular}{|>{\raggedright\arraybackslash}p{0.45\textwidth}|>{\raggedright\arraybackslash}p{0.45\textwidth}|}
\hline
\textbf{Modo atual (busca por texto simples)} & \textbf{Modo desejado (linguagem natural)} \\
\hline
O usuário precisa digitar os termos técnicos que coincidem com os campos do banco, como ``SCN Carta Topográfica Matricial 1:25.000 Rio de Janeiro 2023'' & \textit{``mapas do Rio de Janeiro em 25k publicados em 2023''} \\
\hline
\end{tabular}
\fonte{Adaptado de \citeonline{cgeo2025prototipo}.}
\end{fixo}

Essa barreira técnica restringe o uso do acervo por analistas, gestores e tomadores de decisão que precisam de informação geoespacial, mas não têm formação em cartografia, e problema análogo é relatado na descoberta de dados em geoportais e em catálogos de dados espaciais em geral \cite{hu2015metadata, li2025question}. Por essa razão, no âmbito da Engenharia Cartográfica, o trabalho situa-se na disseminação da informação geoespacial, etapa que torna o acervo produzido acessível a quem dele necessita.

A solução desejável para esse problema é uma busca em linguagem natural apoiada em modelos amplos de linguagem, designados no restante do texto como modelos de linguagem, que interpretam intenções implícitas, normalizam variações de escrita e traduzem descrições informais em estruturas formais \cite{brown2020gpt3, wei2022emergent}. Uma alternativa seria a busca semântica por vetores sobre os metadados, mas critérios como a escala, o código da folha e o intervalo de datas pedem filtros exatos, que a tradução em parâmetros tipados preserva. Nessa busca, o modelo traduz a consulta livre do usuário nos critérios do catálogo, e o sistema executa a busca com esses critérios. Duas das três formulações descritas na Seção~\ref{sec:tc-se-sql} permitem obter dos modelos critérios de busca sem que o modelo escreva a consulta ao banco. No \textit{tool calling}, o modelo recebe a descrição de funções externas, denominadas ferramentas, e escolhe se chama alguma delas e com quais argumentos. O sistema que hospeda o modelo executa a ferramenta chamada e pode devolver o resultado ao modelo antes da resposta final \cite{schick2023toolformer, qin2023toolllm}. Na saída estruturada, o modelo devolve sempre um objeto num formato predefinido pela aplicação, com os campos que preencheu, e o objeto pode ser conferido contra esse formato \cite{tam2024speak, geng2025jsonschemabench}. As duas formulações são detalhadas na Seção~\ref{sec:toolcalling}.

Para tratar o problema, o 1º Centro de Geoinformação (1º CGEO) desenvolveu o projeto \texttt{prototipo\_busca\_llm} \cite{cgeo2025prototipo}, de autoria do Major Diniz. Esse projeto é designado no restante do texto como o protótipo. O protótipo implementa uma busca em linguagem natural sobre o catálogo com o modelo Phi-4 14B \cite{abdin2024phi4} e saída estruturada, com um dicionário de normalização que uniformiza termos da consulta antes de ela chegar ao modelo. O modelo é executado localmente com o Ollama \cite{ollama2023}, plataforma de execução de modelos abertos no próprio \textit{hardware}, e a execução local evita que as consultas trafeguem por serviços externos, condição posta pelo 1º CGEO nas demandas deste trabalho (requisito não funcional RNF1, Seção~\ref{sec:requisitos}), e dispensa o custo recorrente dos serviços em nuvem. O protótipo demonstrou a viabilidade da abordagem no domínio cartográfico nacional e é o ponto de partida deste trabalho. O seu método, reproduzido neste trabalho, constitui o termo de comparação dos resultados. A arquitetura, as contribuições e as limitações do protótipo são analisadas na Seção~\ref{sec:prototipo}.

O \textit{tool calling} tem potencial para aprimorar o protótipo, e a sua adoção consta das demandas que o 1º CGEO apresentou para este trabalho. Com ele, o modelo pode recusar, sem chamar a ferramenta de busca, o pedido que não pertence ao acervo, ao passo que, na saída estruturada, toda consulta resulta num objeto de busca, e a recusa exige um campo criado para esse fim no formato. O modelo pode também chamar ferramentas auxiliares do domínio antes de buscar, como a lista de municípios, o índice de folhas e as regras de datas, e receber o retorno da verificação dos parâmetros como resultado da própria chamada de ferramenta, enquanto, na saída estruturada, esse retorno só chega ao modelo como nova tentativa, depois de emitido o objeto inteiro. Essas vantagens são hipóteses, discutidas nas Seções~\ref{sec:tc-definicao} e~\ref{sec:tc-retorno}, que precisam ser testadas contra a saída estruturada nos mesmos modelos e nas mesmas consultas.

O objeto do trabalho é a tradução da consulta em linguagem natural nos parâmetros de busca do catálogo. A tradução é avaliada contra um gabarito, a resposta correta de cada consulta expressa em parâmetros, e a busca que esses parâmetros produzem é verificada também no acervo real do BDGEx (Seção~\ref{sec:res-acervo}). A busca restringe os resultados pela interseção espacial entre a folha e o polígono do estado, do município (malhas do Instituto Brasileiro de Geografia e Estatística, IBGE) ou da área de suprimento citados na consulta (Seção~\ref{sec:predicados}). Os predicados espaciais relacionais, como ``perto de'', ``num raio de'' e ``ao norte de'', exigiriam geocodificar o lugar citado, isto é, obter a sua posição a partir do nome, além de novos parâmetros e de operadores de distância e de direção, e por isso são excluídos do escopo e indicados como trabalho futuro (Seção~\ref{sec:trabalhos-futuros}). Os doze parâmetros de busca do protótipo foram mantidos para que a solução desenvolvida e o método do protótipo fossem avaliados contra o mesmo gabarito. Também estão fora do escopo:
\begin{itemize}
  \item a interface gráfica para o usuário final;
  \item a implantação em produção na DSG;
  \item a avaliação com usuários e a pesquisa de opinião, previstas nas demandas do 1º CGEO e indicadas como trabalho futuro;
  \item o treinamento e o ajuste fino de modelos de linguagem;
  \item a busca por coordenadas sem referência nominal a um lugar;
  \item idiomas além do português.
\end{itemize}

% ---
\section{Objetivo}
\label{sec:objetivo}
% ---

\subsection{Objetivo Geral}
\label{sec:objetivo-geral}

O objetivo geral do trabalho é adequar o protótipo do 1º CGEO para empregar os recursos de \textit{tool calling} na tradução de consultas em linguagem natural para critérios de busca no acervo cartográfico da DSG, com a finalidade de facilitar a disseminação desse acervo a usuários não especialistas.

\subsection{Objetivos Específicos}
\label{sec:objetivos-especificos}

Para atingir o objetivo geral, foram definidos os seguintes objetivos específicos:
\begin{enumerate}
  \item reescrever o \textit{backend} do protótipo, a parte do sistema que recebe a consulta, aciona o modelo e consulta o banco, para empregar o \textit{tool calling} nativo, no qual o servidor de modelos devolve as chamadas de ferramenta num campo próprio da resposta (Seção~\ref{sec:tc-definicao});
  \item modelar as funcionalidades da busca como ferramentas do domínio, que codificam as convenções da cartografia sistemática e do catálogo;
  \item avaliar modelos de linguagem abertos executados localmente quanto à acurácia, definida como a proporção de consultas traduzidas corretamente, e quanto ao tempo de processamento, e comparar a solução desenvolvida com o protótipo e, num teste A/B, com a saída estruturada executada nos mesmos modelos e nas mesmas consultas;
  \item disponibilizar o código, os conjuntos de consultas construídos para a avaliação e os registros das avaliações.
\end{enumerate}

\subsection{Questões do Trabalho}
\label{sec:questoes}

A avaliação foi planejada para responder às seguintes questões:
\begin{enumerate}
  \item O \textit{tool calling} melhora a tradução em relação ao protótipo?
  \item No teste A/B, o \textit{tool calling} obtém acurácia maior que a da saída estruturada?
  \item No teste A/B, o \textit{tool calling} reduz o tempo de processamento em relação à saída estruturada?
  \item As ferramentas auxiliares do domínio aumentam a acurácia dos modelos locais?
\end{enumerate}

As respostas estão na Seção~\ref{sec:respostas}, na mesma ordem. A tradução foi desenvolvida em três versões sucessivas (Seção~\ref{sec:visao-versoes}), e as respostas às questões 1 e 2 distinguem a primeira versão, a segunda e a versão final. Como análise complementar, sem novas rodadas de modelos, verificou-se também se as buscas produzidas devolvem no acervo real do BDGEx as mesmas cartas que as buscas do gabarito (Seção~\ref{sec:res-acervo}). Essa análise, que não integra as questões planejadas do trabalho, segue o princípio da avaliação pela execução (Seção~\ref{sec:avaliacao-conceitos}), no qual se compara o resultado da busca, e não os parâmetros que a produziram.

% ---
\section{Estrutura do Texto}
\label{sec:estrutura}
% ---

O texto está organizado em cinco capítulos. O Capítulo~\ref{cap:fundamentacao} apresenta os modelos de linguagem empregados, a análise do protótipo e o \textit{tool calling}, com os termos usados na metodologia e na avaliação. O Capítulo~\ref{cap:desenvolvimento} define os requisitos, o ambiente computacional e a sequência metodológica, com o uso de ferramentas de inteligência artificial em cada fase. O Capítulo~\ref{cap:resultados} descreve os dados extraídos do acervo, as funcionalidades da busca e o raciocínio que motivou cada ferramenta, a instalação, a validação e a avaliação dos modelos, a discussão dos resultados e os entregáveis. O Capítulo~\ref{cap:conclusao} retoma o objetivo e apresenta as contribuições, as limitações e os trabalhos futuros.
